\section{Experiments and Results}
\label{sec:experiments}

\subsection{Experimental setup}
Each cohort contains matched histopathology features, gene-expression features, and labels (Table~\ref{tab:datasets}). Our primary experiments use \textbf{TCGA-BRCA} \cite{weinstein2013tcga}, including 113 patients with both modalities. Histopathology is represented by slide-mean-pooled \textbf{UNI2-h} features \cite{chen2024uni}, and genomics by the 200 highest-variance RNA-seq genes. The TCGA-BRCA task predicts \textbf{estrogen-receptor (ER) status}, a molecular
binary endpoint with ER-negative as the positive class ($\approx21\%$), for which gene-only achieves an AUC of $0.95$. We further evaluate generalization on \textbf{TCGA-LUAD} using TP53 mutation prediction in 460 patients, with 48\% positives, and stage classification in 393 patients, distinguishing early stage from stage II or higher. Genomic features are highly predictive for TP53 mutation but less informative for stage. All splits are patient-level. Within each fold, normalization, top-variance
gene selection, PCA, and gene-target construction are fitted using training
patients only.

\begin{table}[t]
\centering
\caption{Paired TCGA cohorts using slide-mean-pooled UNI2-h features and the
200 highest-variance RNA-seq genes.}
\label{tab:datasets}
\resizebox{\linewidth}{!}{%
\begin{tabular}{llccccc}
\toprule
Cohort & Endpoint & Patients & Image dim & Gene dim & Classes & Pos.\ rate \\
\midrule
TCGA-BRCA (UNI2-h) & ER status & 113 & 1536 & 200 & 2 & 21\% \\
TCGA-LUAD (UNI2-h) & TP53 mut. & 460 & 1536 & 200 & 2 & 48\% \\
TCGA-LUAD (UNI2-h) & stage I vs II$+$ & 393 & 1536 & 200 & 2 & 44\% \\
\bottomrule
\end{tabular}%
}
\end{table}


The encoders produce $d{=}128$-dimensional latents. The diffusion predictor uses
FiLM conditioning and a cosine schedule ($T{=}1000$). Each forward pass draws
$N{=}8$ DDIM samples \cite{song2021ddim}, each with eight denoising steps per
direction, to estimate the posterior mean and uncertainty. EMA teachers use decay $0.996$ and
provide stop-gradient targets, and training uses AdamW on a single GPU. A
deterministic MLP attains a comparable AUC, indicating that diffusion is not
required for point prediction alone. Its role is instead to provide sample-based
patient uncertainty for CoRE and to represent the one-to-many
histology-to-genomics mapping.

\paragraph{Baselines and evaluation.}
We compare image-only and gene-only models with early, late, concatenation,
gated, and cross-attention fusion, contrastive alignment, and TabTransformer
\cite{huang2020tabtransformer}, using identical encoders and training settings.
We report AUC, F1, ECE, Brier score, and NLL. Results are mean$\pm$std over five
seeds, with Welch $t$-tests on per-seed values. ECE uses 15-bin top-label
calibration \cite{guo2017calibration} on pooled out-of-fold predictions without
post-hoc scaling; because ECE is binning-sensitive on the small BRCA cohort, we
treat the binning-free Brier score and NLL as the primary calibration evidence.
